\chapter{Event Procedures}
During the execution of CFS++ there are defined points ({\em events})
when interaction between CFS++ and Tcl is possible. These points are
defined by event procedures, which will be called in the .tcl
file. For example, after initialization of the PDEs, the method
\begin{verbatim}
proc CFS_Init {problemName} {
  # put own code here
  # .....
}
\end{verbatim}
will be called, with the name of the problem as input parameter. 
\newline
In CFS there are the following event procedures which are called
during program execution:

\begin{center}
\begin{longtable}{llp{0.5\textwidth}}
\hline\noalign{\smallskip}
procedure & parameters & meaning\\
\noalign{\smallskip}\hline\noalign{\smallskip}
\verb#CFS_Init# & \verb#problemName# (string) 
                & Name of the problem, i.e. basename
                 of the .xml file. \\
\noalign{\smallskip}\hline\noalign{\smallskip}
\verb#CFS_PdeInit# & \verb#pdeName# (string)  
                   & Name of the PDE (e.g. mechanic,
                   electrostatic, direct, acoustic, etc.). Corresponds to the
name
                   written in .log-file.\\
\noalign{\smallskip}\hline\noalign{\smallskip}
\verb#CFS_ReadBCs# & \verb#pdeName# (string)  
                   & Name of the PDE (e.g. mechanic,
                   electrostatic, direct, acoustic, etc.). Corresponds to the name
                   written in .log-file.\\
\noalign{\smallskip}\hline
\verb#CFS_SetBCs# & \verb#pdeName# (string)  
                  & Name of the PDE (e.g. mechanic,
                  electrostatic, direct, acoustic, etc.). Corresponds to the name
                  written in .log-file.\\
                  & \verb#actStep# (unsigned int) 
                  & Current time/frequency step of simulation. \\
                  & \verb#time/frequency# (double)  
                  & Current time/frequency  (dependent on analysis ) \\
\noalign{\smallskip}\hline
\verb#CFS_AssembleMat# & \verb#formName# (string)  
                  & Name of the bilinearform.\\
                  & \verb#elemNum1# (unsigned int) 
                  & First element number. \\
                  & \verb#elemNum2# (unsigned int) 
                  & Second element number. \\
                  & \verb#elemMat# (string)  
                  & Element matrix (pretty printed). \\
\noalign{\smallskip}\hline
\verb#CFS_AssembleRhs# & \verb#formName# (string)  
                  & Name of the linearform.\\
                  & \verb#elemNum# (unsigned int) 
                  & Element number. \\
                  & \verb#elemVec# (string)  
                  & Element vector (pretty printed). \\
\noalign{\smallskip}\hline
\verb#CFS_PostProcess# & \verb#pdeName# (string)  
                       & Name of the PDE (e.g. mechanic,
                       electrostatic, direct, acoustic, etc.). Corresponds to the name
                       written in .log-file.\\
\caption{Overview of Tcl event procedures}
\label{tab:tcl-events}
\end{longtable}

\end{center}

Some notes about the event procedures:
\begin{itemize}
\item By default, all event procedures are defined as empty
functions. Only the ones overwritten / defined in the .tcl file will
actually perform any {\em real} operation.
\item Be careful to specify the correct number of parameters as given
  in table \ref{tab:tcl-events} as otherwise the method will not be
  recognized as event procedure!
\item Event procedures will generally be called {\bf after} the
  parameters from the .xml file are read. For example, the method
  {\em CFS\_ReadBCs} will be called after the parameters were read from
  the .xml file.
\end{itemize}